\section{Vanishing of cohomology}
\label{section-vanishing}

\noindent
We have seen that on an affine scheme the higher cohomology groups
of any quasi-coherent sheaf vanish
(Lemma \ref{lemma-quasi-coherent-affine-cohomology-zero}).
It turns out that this also
characterizes affine schemes. We give two versions.

\begin{lemma}
\label{lemma-quasi-compact-h1-zero-covering}
\begin{reference}
\cite{Serre-criterion}, \cite[II, Theorem 5.2.1 (d') and IV (1.7.17)]{EGA}
\end{reference}
\begin{slogan}
Serre's criterion for affineness.
\end{slogan}
Let $X$ be a scheme.
Assume that
\begin{enumerate}
\item $X$ is quasi-compact,
\item for every quasi-coherent sheaf of ideals
$\mathcal{I} \subset \mathcal{O}_X$ we have $H^1(X, \mathcal{I}) = 0$.
\end{enumerate}
Then $X$ is affine.
\end{lemma}

\begin{proof}
Let $x \in X$ be a closed point. Let $U \subset X$ be an affine open
neighbourhood of $x$. Write $U = \Spec(A)$ and let
$\mathfrak m \subset A$ be the maximal ideal corresponding to $x$.
Set $Z = X \setminus U$ and $Z' = Z \cup \{x\}$.
By Schemes, Lemma \ref{schemes-lemma-reduced-closed-subscheme} there
are quasi-coherent sheaves of ideals
$\mathcal{I}$, resp.\ $\mathcal{I}'$ cutting out
the reduced closed subschemes $Z$, resp.\ $Z'$.
Consider the short exact sequence
$$
0 \to \mathcal{I}' \to \mathcal{I} \to \mathcal{I}/\mathcal{I}' \to 0.
$$
Since $x$ is a closed point of $X$ and $x \not \in Z$ we see that
$\mathcal{I}/\mathcal{I}'$ is supported at $x$. In fact, the restriction
of $\mathcal{I}/\mathcal{I'}$ to $U$ corresponds to the $A$-module
$A/\mathfrak m$. Hence we see that $\Gamma(X, \mathcal{I}/\mathcal{I'})
= A/\mathfrak m$. Since by assumption $H^1(X, \mathcal{I}') = 0$
we see there exists a global section $f \in \Gamma(X, \mathcal{I})$
which maps to the element $1 \in A/\mathfrak m$ as a section of
$\mathcal{I}/\mathcal{I'}$. Clearly we have
$x \in X_f \subset U$. This implies that $X_f = D(f_A)$ where
$f_A$ is the image of $f$ in $A = \Gamma(U, \mathcal{O}_X)$.
In particular $X_f$ is affine.

\medskip\noindent
Consider the union $W = \bigcup X_f$ over all $f \in \Gamma(X, \mathcal{O}_X)$
such that $X_f$ is affine. Obviously $W$ is open in $X$.
By the arguments above every closed point of
$X$ is contained in $W$. The closed subset $X \setminus W$ of $X$
is also quasi-compact
(see Topology, Lemma \ref{topology-lemma-closed-in-quasi-compact}).
Hence it has a closed point if it is nonempty (see
Topology, Lemma \ref{topology-lemma-quasi-compact-closed-point}).
This would contradict the fact that all closed points are in
$W$. Hence we conclude $X = W$.

\medskip\noindent
Choose finitely many $f_1, \ldots, f_n \in \Gamma(X, \mathcal{O}_X)$
such that $X = X_{f_1} \cup \ldots \cup X_{f_n}$ and such that each
$X_{f_i}$ is affine. This is possible as we've seen above.
By Properties, Lemma \ref{properties-lemma-characterize-affine}
to finish the proof it suffices
to show that $f_1, \ldots, f_n$ generate the unit ideal in
$\Gamma(X, \mathcal{O}_X)$. Consider the short exact sequence
$$
\xymatrix{
0 \ar[r] &
\mathcal{F} \ar[r] &
\mathcal{O}_X^{\oplus n} \ar[rr]^{f_1, \ldots, f_n} & &
\mathcal{O}_X \ar[r] &
0
}
$$
The arrow defined by $f_1, \ldots, f_n$ is surjective since the
opens $X_{f_i}$ cover $X$. We let $\mathcal{F}$ be the kernel
of this surjective map.
Observe that $\mathcal{F}$ has a filtration
$$
0 = \mathcal{F}_0 \subset \mathcal{F}_1 \subset
\ldots \subset \mathcal{F}_n = \mathcal{F}
$$
so that each subquotient $\mathcal{F}_i/\mathcal{F}_{i - 1}$ is
isomorphic to a quasi-coherent sheaf of ideals.
Namely we can take $\mathcal{F}_i$ to be the intersection of
$\mathcal{F}$ with the first $i$ direct summands of
$\mathcal{O}_X^{\oplus n}$.
The assumption
of the lemma implies that $H^1(X, \mathcal{F}_i/\mathcal{F}_{i - 1}) = 0$
for all $i$. This implies that
$H^1(X, \mathcal{F}_2) = 0$ because it is sandwiched between
$H^1(X, \mathcal{F}_1)$ and $H^1(X, \mathcal{F}_2/\mathcal{F}_1)$.
Continuing like this we deduce that $H^1(X, \mathcal{F}) = 0$.
Therefore we conclude that the map
$$
\xymatrix{
\bigoplus\nolimits_{i = 1, \ldots, n} \Gamma(X, \mathcal{O}_X)
\ar[rr]^{f_1, \ldots, f_n} & &
\Gamma(X, \mathcal{O}_X)
}
$$
is surjective as desired.
\end{proof}

\noindent
Note that if $X$ is a Noetherian scheme then every quasi-coherent
sheaf of ideals is automatically a coherent sheaf of ideals and a
finite type quasi-coherent sheaf of ideals. Hence
the preceding lemma and the next lemma both apply in this case.

\begin{lemma}
\label{lemma-quasi-separated-h1-zero-covering}
\begin{reference}
\cite{Serre-criterion}, \cite[II, Theorem 5.2.1]{EGA}
\end{reference}
\begin{slogan}
Serre's criterion for affineness.
\end{slogan}
Let $X$ be a scheme. Assume that
\begin{enumerate}
\item $X$ is quasi-compact,
\item $X$ is quasi-separated, and
\item $H^1(X, \mathcal{I}) = 0$ for every quasi-coherent sheaf
of ideals $\mathcal{I}$ of finite type.
\end{enumerate}
Then $X$ is affine.
\end{lemma}

\begin{proof}
By
Properties, Lemma \ref{properties-lemma-quasi-coherent-colimit-finite-type}
every quasi-coherent sheaf of ideals is a directed colimit of
quasi-coherent sheaves of ideals of finite type.
By Cohomology, Lemma \ref{cohomology-lemma-quasi-separated-cohomology-colimit}
taking cohomology on $X$ commutes with directed colimits.
Hence we see that $H^1(X, \mathcal{I}) = 0$
for every quasi-coherent sheaf of ideals on $X$. In other words
we see that Lemma \ref{lemma-quasi-compact-h1-zero-covering} applies.
\end{proof}

\noindent
We can use the arguments given above to find a sufficient condition to
see when an invertible sheaf is ample. However, we warn the reader that
this condition is not necessary.

\begin{lemma}
\label{lemma-quasi-compact-h1-zero-invertible}
Let $X$ be a scheme. Let $\mathcal{L}$ be an invertible $\mathcal{O}_X$-module.
Assume that
\begin{enumerate}
\item $X$ is quasi-compact,
\item for every quasi-coherent sheaf of ideals
$\mathcal{I} \subset \mathcal{O}_X$
there exists an $n \geq 1$ such that
$H^1(X, \mathcal{I} \otimes_{\mathcal{O}_X} \mathcal{L}^{\otimes n}) = 0$.
\end{enumerate}
Then $\mathcal{L}$ is ample.
\end{lemma}

\begin{proof}
This is proved in exactly the same way as
Lemma \ref{lemma-quasi-compact-h1-zero-covering}.
Let $x \in X$ be a closed point. Let $U \subset X$ be an affine open
neighbourhood of $x$ such that $\mathcal{L}|_U \cong \mathcal{O}_U$.
Write $U = \Spec(A)$ and let
$\mathfrak m \subset A$ be the maximal ideal corresponding to $x$.
Set $Z = X \setminus U$ and $Z' = Z \cup \{x\}$.
By Schemes, Lemma \ref{schemes-lemma-reduced-closed-subscheme} there
are quasi-coherent sheaves of ideals
$\mathcal{I}$, resp.\ $\mathcal{I}'$ cutting out
the reduced closed subschemes $Z$, resp.\ $Z'$.
Consider the short exact sequence
$$
0 \to \mathcal{I}' \to \mathcal{I} \to \mathcal{I}/\mathcal{I}' \to 0.
$$
For every $n \geq 1$ we obtain a short exact sequence
$$
0 \to \mathcal{I}' \otimes_{\mathcal{O}_X} \mathcal{L}^{\otimes n}
\to \mathcal{I} \otimes_{\mathcal{O}_X} \mathcal{L}^{\otimes n} \to
\mathcal{I}/\mathcal{I}' \otimes_{\mathcal{O}_X} \mathcal{L}^{\otimes n} \to 0.
$$
By our assumption we may pick $n$ such that
$H^1(X, \mathcal{I}' \otimes_{\mathcal{O}_X} \mathcal{L}^{\otimes n}) = 0$.
Since $x$ is a closed point of $X$ and $x \not \in Z$ we see that
$\mathcal{I}/\mathcal{I}'$ is supported at $x$. In fact, the restriction
of $\mathcal{I}/\mathcal{I'}$ to $U$ corresponds to the $A$-module
$A/\mathfrak m$. Since $\mathcal{L}$ is trivial on $U$
we see that the restriction of
$\mathcal{I}/\mathcal{I}' \otimes_{\mathcal{O}_X} \mathcal{L}^{\otimes n}$
to $U$ also corresponds to the $A$-module $A/\mathfrak m$.
Hence we see that
$\Gamma(X, \mathcal{I}/\mathcal{I'} \otimes_{\mathcal{O}_X}
\mathcal{L}^{\otimes n}) = A/\mathfrak m$.
By our choice of $n$ we see there exists a global section
$s \in \Gamma(X, \mathcal{I} \otimes_{\mathcal{O}_X} \mathcal{L}^{\otimes n})$
which maps to the element $1 \in A/\mathfrak m$. Clearly we have
$x \in X_s \subset U$ because $s$ vanishes at points of $Z$.
This implies that $X_s = D(f)$ where
$f \in A$ is the image of $s$ in $A \cong \Gamma(U, \mathcal{L}^{\otimes n})$.
In particular $X_s$ is affine.

\medskip\noindent
Consider the union $W = \bigcup X_s$ over all
$s \in \Gamma(X, \mathcal{L}^{\otimes n})$ for $n \geq 1$
such that $X_s$ is affine. Obviously $W$ is open in $X$.
By the arguments above every closed point of
$X$ is contained in $W$. The closed subset $X \setminus W$ of $X$
is also quasi-compact
(see Topology, Lemma \ref{topology-lemma-closed-in-quasi-compact}).
Hence it has a closed point if it is nonempty (see
Topology, Lemma \ref{topology-lemma-quasi-compact-closed-point}).
This would contradict the fact that all closed points are in
$W$. Hence we conclude $X = W$. This means that $\mathcal{L}$
is ample by Properties, Definition \ref{properties-definition-ample}.
\end{proof}

\noindent
There is a variant of Lemma \ref{lemma-quasi-compact-h1-zero-invertible}
with finite type ideal sheaves which we will formulate and prove here if
we ever need it.

\begin{lemma}
\label{lemma-criterion-affine-morphism}
Let $f : X \to Y$ be a quasi-compact morphism with $X$ and $Y$ quasi-separated.
If $R^1f_*\mathcal{I} = 0$ for every quasi-coherent sheaf of ideals
$\mathcal{I}$ on $X$, then $f$ is affine.
\end{lemma}

\begin{proof}
Let $V \subset Y$ be an affine open subscheme. We have to show that
$U = f^{-1}(V)$ is affine. The inclusion morphism $V \to Y$ is quasi-compact
by Schemes, Lemma \ref{schemes-lemma-quasi-compact-permanence}.
Hence the base change $U \to X$ is quasi-compact, see
Schemes, Lemma \ref{schemes-lemma-quasi-compact-preserved-base-change}.
Thus any quasi-coherent sheaf of ideals $\mathcal{I}$ on $U$
extends to a quasi-coherent sheaf of ideals on $X$, see
Properties, Lemma \ref{properties-lemma-extend-trivial}.
Since the formation of $R^1f_*$ is local on $Y$
(Cohomology, Section \ref{cohomology-section-locality})
we conclude that $R^1(U \to V)_*\mathcal{I} = 0$ by the assumption
in the lemma. Hence by the Leray Spectral sequence
(Cohomology, Lemma \ref{cohomology-lemma-Leray})
we conclude that $H^1(U, \mathcal{I}) = H^1(V, (U \to V)_*\mathcal{I})$.
Since $(U \to V)_*\mathcal{I}$ is quasi-coherent by
Schemes, Lemma \ref{schemes-lemma-push-forward-quasi-coherent}, we have
$H^1(V, (U \to V)_*\mathcal{I}) = 0$ by
Lemma \ref{lemma-quasi-coherent-affine-cohomology-zero}.
Thus we find that $U$ is affine by
Lemma \ref{lemma-quasi-compact-h1-zero-covering}.
\end{proof}















